\chapter{Conclusions}
\label{ch:thisChapterLabel}
\section{Higher-Order Reconstruction Results}

The Roe scheme adds a large amount of diffusion in regions near where rarefactions and contact discontinuities occur. In all of the test cases, the Roe-WENO5 scheme provides the least diffusive solution. The improvement provided by the Roe-WENO5 scheme is most apparent in regions containing contact discontinuities. Near contact discontinuities, the Roe-WENO5 scheme can be seen to add much less diffusion than the Roe, Roe-MUSCL, and Roe-WENO3 schemes. This reduction in numerical diffusion leads to better capturing of the contact wave. The exception to this improvement is the standing contact wave in Test 4, where the Roe-WENO5 scheme only provides slightly less diffusion than the Roe, Roe-MUSCL, and Roe-WENO3 schemes.

The Roe-WENO5 scheme adds less numerical diffusion near rarefactions than the Roe, Roe-MUSCL, and Roe-WENO3. This leads to a more accurate capturing of rarefactions than the other three schemes. The Roe scheme adds the most diffusion near rarefactions, leading to inaccurate capturing of the rarefaction wave. The Roe-WENO5 adds the least diffusion near the rarefaction waves and captures the wave much more accurately. The Roe-WENO3 scheme captures the rarefaction wave more accurately than the Roe and Roe-MUSCL schemes, but the Roe-WENO5 scheme can be seen to be an improvement.

In cases such as Test 2, where the contact wave is very close to a shock, numerical methods may not accurately represent the region of constant density between the contact wave and shock wave due to a large amount of diffusion added near the contact wave. In this case, the Roe-WENO5 scheme adds much less diffusion near the contact wave and provides a much more accurate solution. In Test 4, however, this is not the case. The Roe, Roe-MUSCL, and Roe-WENO3 schemes also provide an accurate solution in this region. The Roe-WENO5 scheme also displays some small pressure oscillations near the location of contact discontinuities for the Roe-MUSCL scheme and absent from the Roe and Roe-MUSCL schemes.

The Roe-WENO5 scheme captures shocks better than the Roe, Roe-MUSCL, and Roe-WENO3 schemes; however, the advantage over the Roe-MUSCL and Roe-WENO3 schemes are minimal. The Roe scheme leads to a smearing of the solution in the region. This smearing is easily observed in Test 1, Test 2, and Test 3. The Roe-MUSCL and Roe-WENO3 schemes can be seen to capture shock waves equally well, adding only slightly more diffusion near shocks than the Roe-WENO5 scheme.
\section{Higher-Order Reconstruction with Higher-Order Central Differencing Results}

The Roe scheme was the least sensitive to the modification made to the average flux term. The results for the CD2-Roe scheme remained unchanged from the results for the Roe scheme in all tests. The average flux in the Roe scheme is seen in \equationref{eqn:CDFlux} where $\boldsymbol{F}_L$ is taken as $\boldsymbol{F}_{i-1}$ and $\boldsymbol{F}_R$ is taken as $\boldsymbol{F}_i$. The centrally differenced flux in the CD2 scheme is seen in \equationref{eqn:CD2}, which is equivalent to the average flux in the Roe scheme. 

The CD2-Roe-WENO5 scheme shows large oscillations in all four tests, with the largest amplitude oscillations near contact discontinuities. The CD2-Roe-MUSCL and CD2-Roe-WENO3 schemes are not affected to the same degree as the CD2-Roe-WENO5 scheme. The CD2-Roe-MUSCL and CD2-Roe-WENO3 schemes show some small oscillatory behavior near shocks in Test 3, but the oscillations are of much lower amplitude than seen in the CD2-Roe-WENO5 scheme. Near contact waves, the CD2-Roe-MUSCL and CD2-Roe-WENO3 schemes also show smaller oscillations than the CD2-Roe-WENO5 scheme. The CD4-Roe-WENO5 scheme sees the oscillations near the contact waves dampen when compared to the CD2-Roe-WENO5 scheme. Test 3, however, also displays an increase in the oscillations' amplitude ahead of the left shock. The oscillations near contact waves in the CD6-Roe-WENO5 scheme can also be seen to dampen when compared to the CD4-Roe-WENO5 scheme. The oscillations near the left shock in Test 3 increase in amplitude as well.

In Test 1, the higher-order reconstruction with higher-order central differencing schemes showed the reintroduction of the entropy fix artifact present in the Roe scheme absent from the higher-order reconstruction schemes. Except for the CD2-Roe scheme, which is, in fact, the Roe scheme, the CD2, CD4, and CD6 schemes did not produce solutions for Test 2 and Test 4. These schemes developed large oscillations at the contact waves that led to non-TVD behavior and unphysical flow properties, making them unsuitable for use in any practical computation.
\section{Conclusions}

The Roe-WENO5 consistently produces the least diffusive and most accurate solutions of the higher-order reconstruction schemes. The Roe-WENO3 scheme displays less oscillatory behavior than the Roe-WENO5 scheme; however, the Roe-WENO5 scheme captures contact discontinuities more accurately than the Roe-WENO3 scheme.  The Roe scheme adds too much diffusion near rarefactions and contact waves to accurately capture these regions without an increase in resolution. Increasing the number of grid points in the computational domain to resolve these flow features leads to costly increases in computation time.

The Roe-WENO5 scheme provides the lowest numerical diffusion while keeping the solution TVD of the schemes tested. The higher-order central differencing schemes proved too unstable for use in practical computations. The CD2, CD4, and CD6 with higher-order reconstruction schemes had large oscillations near contact discontinuities that led to the schemes' failure in Test 2 and Test 4, and in Test 3 developed large oscillations near the left shock wave. In future work, another method of increasing the order of the average flux is needed. The approach of increasing the order using higher-order central differencing has proven to be prone to oscillations that can cause the solution to lose accuracy or fail.